\documentclass[12pt, a4paper]{article}

% ── Packages ────────────────────────────────────────────────────────────────
\usepackage[utf8]{inputenc}
\usepackage[T1]{fontenc}
\usepackage[margin=1in]{geometry}
\usepackage{xcolor}
\usepackage{hyperref}
\usepackage{listings}
\usepackage{amsmath}
\usepackage{booktabs}
\usepackage{fancyhdr}
\usepackage{titlesec}
\usepackage{parskip}
\usepackage{graphicx}

% ── Colours ─────────────────────────────────────────────────────────────────
\definecolor{codebackground}{HTML}{F5F5F5}
\definecolor{codekeyword}{HTML}{0060A8}
\definecolor{codecomment}{HTML}{6A737D}
\definecolor{codestring}{HTML}{B31D28}
\definecolor{codenumber}{HTML}{005CC5}
\definecolor{accentblue}{HTML}{0366D6}

% ── Code listing style ───────────────────────────────────────────────────────
\lstdefinestyle{python}{
    language=Python,
    basicstyle=\ttfamily\small,
    keywordstyle=\color{codekeyword}\bfseries,
    commentstyle=\color{codecomment}\itshape,
    stringstyle=\color{codestring},
    numberstyle=\tiny\color{codenumber},
    backgroundcolor=\color{codebackground},
    breaklines=true,
    numbers=left,
    numbersep=8pt,
    frame=single,
    rulecolor=\color{black!20},
    tabsize=4,
    showstringspaces=false,
    captionpos=b,
    xleftmargin=2em,
    framexleftmargin=1.5em,
    aboveskip=0.8em,
    belowskip=0.5em,
    morekeywords={yield,from,lambda,frozenset,namedtuple},
}
\lstset{style=python}

% ── Python Tutor placeholder box ─────────────────────────────────────────────
\newcommand{\pythontutorbox}{%
  \vspace{0.5em}%
  \noindent\fbox{\parbox{\linewidth}{%
    \centering%
    \vspace{4cm}%
    \textit{\color{gray} [ Python Tutor Visualization --- paste the code above at \texttt{https://pythontutor.com} ]}%
    \vspace{4cm}%
  }}%
  \vspace{0.8em}%
}

% ── Header / Footer ─────────────────────────────────────────────────────────
\pagestyle{fancy}
\fancyhf{}
\fancyhead[L]{\small\textit{Principles of Functional Languages}}
\fancyhead[R]{\small\thepage}
\renewcommand{\headrulewidth}{0.4pt}

% ── Section formatting ───────────────────────────────────────────────────────
\titleformat{\section}{\large\bfseries}{}{0em}{}[\vspace{-0.3em}\rule{\linewidth}{0.4pt}]

\hypersetup{colorlinks=true, linkcolor=accentblue, urlcolor=accentblue}

% ════════════════════════════════════════════════════════════════════════════
\begin{document}

% ── Title Page ───────────────────────────────────────────────────────────────
\begin{titlepage}
\centering
\vspace*{2.5cm}

{\Large\bfseries PRINCIPLES OF FUNCTIONAL LANGUAGES\\[0.5em]
PYTHON TUTOR ASSESSMENT\par}

\vspace{3cm}

{\normalsize
\begin{tabular}{rl}
\textbf{Name:}    & Your Name Here \\[0.3em]
\textbf{Roll No:} & CB.SC.U4CSExxxxx \\
\end{tabular}
\par}

\vspace{3.5cm}

{\Large\bfseries FUNCTIONAL EQUATION PARSER \& ANALYZER\par}

\vfill
{\small\itshape Department of Computer Science\\
Amrita School of Computing}

\end{titlepage}

\tableofcontents
\newpage

% ════════════════════════════════════════════════════════════════════════════
%  1. FUNCTIONS AS OBJECTS
% ════════════════════════════════════════════════════════════════════════════
\section{1.\quad Functions as Objects}

\textbf{Code:}

\begin{lstlisting}
# Functions are first-class objects — assign them to variables
from collections import namedtuple

Number   = namedtuple("Number",   ["value"])
Variable = namedtuple("Variable", ["name"])
BinOp    = namedtuple("BinOp",    ["op", "left", "right"])

# A function stored in a variable
def str_number(node):   return str(node.value)
def str_variable(node): return node.name
def str_binop(node):    return f"({to_str(node.left)} {node.op} {to_str(node.right)})"

# Each function is assigned as a VALUE in a dictionary
TO_STRING = {
    Number:   str_number,
    Variable: str_variable,
    BinOp:    str_binop,
}

def to_str(node):
    handler = TO_STRING[type(node)]   # retrieve function object
    return handler(node)              # call it

tree = BinOp("+", Number(2), Variable("x"))
print(to_str(tree))   # ((2 + x))
\end{lstlisting}

\pythontutorbox

\textbf{Code Explanation:}

In this code, each formatting function (\texttt{str\_number}, \texttt{str\_variable}, \texttt{str\_binop}) is stored as a \textbf{value} inside the dictionary \texttt{TO\_STRING}. This is the \textbf{dispatch-table} pattern.

\textbf{What happens here?}
\begin{itemize}
    \item The three functions are stored as objects (values) in the dictionary.
    \item \texttt{TO\_STRING[type(node)]} retrieves the correct function for the node type.
    \item \texttt{handler(node)} calls it --- treating the function exactly like any other object.
    \item This replaces \texttt{if/elif} chains with a clean, extensible table.
\end{itemize}

\newpage
% ════════════════════════════════════════════════════════════════════════════
%  2. HIGHER-ORDER FUNCTIONS
% ════════════════════════════════════════════════════════════════════════════
\section{2.\quad Higher-Order Functions}

\textbf{Code:}

\begin{lstlisting}
# Higher-order function: takes operators and ANOTHER function as arguments
from collections import namedtuple

BinOp = namedtuple("BinOp", ["op", "left", "right"])
Number = namedtuple("Number", ["value"])
Token = namedtuple("Token", ["kind", "value"])
OP, EOF = "OP", "EOF"

def parse_binary(operators, next_parser):
    """Factory: returns a NEW parser for left-associative op at one precedence level."""
    def _parser(peek, advance):
        left = next_parser(peek, advance)
        while peek().kind == OP and peek().value in operators:
            op = advance().value
            right = next_parser(peek, advance)
            left = BinOp(op, left, right)
        return left
    return _parser   # returns a function!

# Demonstrate: build three parsers from the SAME factory
def dummy_next(peek, advance):
    t = advance()
    return Number(t.value)

parse_sum   = parse_binary({"+", "-"}, dummy_next)
parse_term  = parse_binary({"*", "/"}, dummy_next)
parse_power = parse_binary({"^"},      dummy_next)

print("parse_sum   is a function:", callable(parse_sum))
print("parse_term  is a function:", callable(parse_term))
print("parse_power is a function:", callable(parse_power))
\end{lstlisting}

\pythontutorbox

\textbf{Code Explanation:}

\texttt{parse\_binary} is a \textbf{higher-order function} --- it takes \texttt{next\_parser} (a function) as an argument and \textbf{returns a new function}.

\textbf{What happens here?}
\begin{itemize}
    \item The same factory produces three distinct parsers (\texttt{parse\_sum}, \texttt{parse\_term}, \texttt{parse\_power}).
    \item Adding a new precedence level (e.g.\ modulo \texttt{\%}) requires only one line: \texttt{parse\_mod = parse\_binary(\{"\%"\}, parse\_power)}.
    \item This eliminates repetition while encoding all operator-precedence rules.
\end{itemize}

\newpage
% ════════════════════════════════════════════════════════════════════════════
%  3. IMMUTABILITY
% ════════════════════════════════════════════════════════════════════════════
\section{3.\quad Immutability}

\textbf{Code:}

\begin{lstlisting}
from collections import namedtuple

# All AST nodes are namedtuples -- immutable by design
Number   = namedtuple("Number",   ["value"])
Variable = namedtuple("Variable", ["name"])
BinOp    = namedtuple("BinOp",    ["op", "left", "right"])

# Create a node
node = BinOp("+", Number(2), Variable("x"))
print("Original:", node)

# Attempt to modify -- this will raise AttributeError
try:
    node.op = "-"
except AttributeError as e:
    print("Cannot modify:", e)

# To 'change' a node you must create a NEW one
new_node = BinOp("-", node.left, node.right)
print("New node:", new_node)
print("Original unchanged:", node)
\end{lstlisting}

\pythontutorbox

\textbf{Code Explanation:}

AST nodes are \texttt{namedtuple} instances. A \texttt{namedtuple} is \textbf{immutable} --- its fields cannot be changed after creation, just like a Python \texttt{tuple}.

\textbf{What happens here?}
\begin{itemize}
    \item \texttt{BinOp("+", Number(2), Variable("x"))} creates an immutable node.
    \item Trying \texttt{node.op = "-"} raises \texttt{AttributeError}.
    \item Transformations (simplification, differentiation) always create \textbf{new} nodes rather than mutating existing ones.
\end{itemize}

\newpage
% ════════════════════════════════════════════════════════════════════════════
%  4. GENERATOR
% ════════════════════════════════════════════════════════════════════════════
\section{4.\quad Generator}

\textbf{Code:}

\begin{lstlisting}
# The tokenizer is a generator -- produces tokens one at a time with yield
from collections import namedtuple

Token = namedtuple("Token", ["kind", "value"])
NUMBER, IDENT, OP, EOF = "NUMBER", "IDENT", "OP", "EOF"

def tokenize(text):
    """Generator: yields Token namedtuples lazily."""
    pos = 0
    while pos < len(text):
        ch = text[pos]
        if ch.isspace():
            pos += 1; continue
        if ch.isdigit():
            start = pos
            while pos < len(text) and text[pos].isdigit(): pos += 1
            yield Token(NUMBER, int(text[start:pos]))
            continue
        if ch.isalpha():
            yield Token(IDENT, ch); pos += 1; continue
        if ch in "+-*/^":
            yield Token(OP, ch); pos += 1; continue
        pos += 1
    yield Token(EOF, None)   # sentinel

# Prove it is a generator
gen = tokenize("2 + x")
print("Is generator?", hasattr(gen, "__next__"))

# Tokens produced lazily
t1 = next(gen); print("Token 1:", t1)
t2 = next(gen); print("Token 2:", t2)
t3 = next(gen); print("Token 3:", t3)

# Or iterate over it
for tok in tokenize("3 * y"):
    print(tok)
\end{lstlisting}

\pythontutorbox

\textbf{Code Explanation:}

A \textbf{generator} is a function that uses \texttt{yield} instead of \texttt{return}. It produces values \textit{lazily} --- one at a time, only when requested.

\textbf{What happens here?}
\begin{itemize}
    \item \texttt{tokenize("2 + x")} does \textbf{not} scan the whole string immediately.
    \item Each call to \texttt{next(gen)} resumes the function from where it last yielded.
    \item This is called \textbf{lazy evaluation} --- memory-efficient, processes one token at a time.
\end{itemize}

\newpage
% ════════════════════════════════════════════════════════════════════════════
%  5. ITERATOR
% ════════════════════════════════════════════════════════════════════════════
\section{5.\quad Iterator}

\textbf{Code:}

\begin{lstlisting}
from collections import namedtuple

Token = namedtuple("Token", ["kind", "value"])
EOF = "EOF"

def tokenize(text):
    pos = 0
    while pos < len(text):
        ch = text[pos]
        if ch.isspace():   pos += 1; continue
        if ch.isdigit():
            yield Token("NUMBER", int(ch)); pos += 1; continue
        if ch in "+-*/":
            yield Token("OP", ch); pos += 1; continue
        pos += 1
    yield Token(EOF, None)

# Build a closure-based peekable iterator (no class needed!)
def make_peekable(token_iter):
    buffer = [next(token_iter)]      # mutable state captured by closures

    def peek():                      # look at current token
        return buffer[0]

    def advance():                   # consume and move to next
        current = buffer[0]
        try:
            buffer[0] = next(token_iter)
        except StopIteration:
            buffer[0] = Token(EOF, None)
        return current

    return peek, advance             # return a pair of closures

# Demo
peek, advance = make_peekable(tokenize("5 + 3"))
print("peek:", peek())       # see current without consuming
t = advance()                # consume it
print("advanced:", t)
print("peek after:", peek()) # now at next token
\end{lstlisting}

\pythontutorbox

\textbf{Code Explanation:}

An \textbf{iterator} allows sequential one-at-a-time access to elements. Here, a \textbf{closure pair} (\texttt{peek}, \texttt{advance}) wraps the token generator into a peekable iterator --- without writing a class.

\textbf{What happens here?}
\begin{itemize}
    \item \texttt{make\_peekable} captures mutable state in \texttt{buffer} --- a list with one element.
    \item \texttt{peek()} reads the current token without advancing.
    \item \texttt{advance()} consumes it and loads the next one.
    \item The parser uses this to look one token ahead --- the essential lookahead needed for recursive descent parsing.
\end{itemize}

\newpage
% ════════════════════════════════════════════════════════════════════════════
%  6. CLOSURE
% ════════════════════════════════════════════════════════════════════════════
\section{6.\quad Closure}

\textbf{Code:}

\begin{lstlisting}
# Closures: inner functions that remember their enclosing scope
from collections import namedtuple

Token = namedtuple("Token", ["kind", "value"])
EOF = "EOF"

def make_peekable(token_iter):
    buffer = [next(token_iter)]   # <-- this variable is CAPTURED

    def peek():
        return buffer[0]          # remembers 'buffer' from outer scope

    def advance():
        current = buffer[0]       # remembers 'buffer' from outer scope
        try:
            buffer[0] = next(token_iter)
        except StopIteration:
            buffer[0] = Token(EOF, None)
        return current

    return peek, advance

# Show that the closures share the same 'buffer'
tokens = iter([
    Token("NUMBER", 2),
    Token("OP", "+"),
    Token("NUMBER", 3),
    Token(EOF, None),
])

peek, advance = make_peekable(tokens)
print("Before advance, peek:", peek())   # NUMBER 2
advance()
print("After  advance, peek:", peek())   # OP +
advance()
print("After  advance, peek:", peek())   # NUMBER 3
\end{lstlisting}

\pythontutorbox

\textbf{Code Explanation:}

A \textbf{closure} is an inner function that \textit{remembers} variables from its enclosing function even after the outer function has returned.

\textbf{What happens here?}
\begin{itemize}
    \item \texttt{make\_peekable} returns two functions: \texttt{peek} and \texttt{advance}.
    \item Both functions close over (capture) the \texttt{buffer} variable.
    \item When \texttt{advance()} modifies \texttt{buffer[0]}, \texttt{peek()} immediately sees the new value --- they share the same captured reference.
    \item This pattern replaces a stateful class with two lightweight closure objects.
\end{itemize}

\newpage
% ════════════════════════════════════════════════════════════════════════════
%  7. CURRYING
% ════════════════════════════════════════════════════════════════════════════
\section{7.\quad Currying}

\textbf{Code:}

\begin{lstlisting}
# Currying: transform a two-argument function into a chain of one-arg functions
from functools import reduce

def curry(fn):
    """Two-argument currying decorator."""
    def _outer(a):
        def _inner(b):
            return fn(a, b)      # captures 'a' from outer scope (closure!)
        return _inner
    return _outer

# Example: curry a differentiation dispatcher
@curry
def diff_wrt(var, node_value):
    """Return 1 if differentiating with respect to var, else 0."""
    return 1 if node_value == var else 0

# Create specialised diffentiators by partial application
diff_wrt_x = diff_wrt("x")   # fix first argument to "x"
diff_wrt_y = diff_wrt("y")   # fix first argument to "y"

print(diff_wrt_x("x"))   # 1  -- differentiating x w.r.t. x
print(diff_wrt_x("y"))   # 0  -- differentiating y w.r.t. x
print(diff_wrt_y("y"))   # 1  -- differentiating y w.r.t. y

# Same result calling both args directly
print(diff_wrt("x")("x"))   # 1
\end{lstlisting}

\pythontutorbox

\textbf{Code Explanation:}

\textbf{Currying} converts a function that takes two arguments into a function that takes one argument and \textit{returns} another function expecting the second argument.

\textbf{What happens here?}
\begin{itemize}
    \item \texttt{curry(fn)} wraps \texttt{fn} so it can be called one argument at a time.
    \item \texttt{diff\_wrt("x")} returns a new specialised function \texttt{diff\_wrt\_x}.
    \item \texttt{diff\_wrt\_x("x")} finally calls the original logic.
    \item This enables \textbf{partial application} --- fixing some arguments for reuse.
\end{itemize}

\newpage
% ════════════════════════════════════════════════════════════════════════════
%  8. COMPOSE / PIPE
% ════════════════════════════════════════════════════════════════════════════
\section{8.\quad Function Composition (pipe \& compose)}

\textbf{Code:}

\begin{lstlisting}
from functools import reduce

def compose(*fns):
    """Right-to-left: compose(f, g)(x) == f(g(x))"""
    def _composed(x):
        return reduce(lambda acc, f: f(acc), reversed(fns), x)
    return _composed

def pipe(*fns):
    """Left-to-right: pipe(f, g)(x) == g(f(x))"""
    def _piped(x):
        return reduce(lambda acc, f: f(acc), fns, x)
    return _piped

# Stand-alone functions for the equation pipeline
def parse(text):
    # simplified: just wrap in a dict for demo
    return {"type": "expr", "raw": text}

def simplify(node):
    node["simplified"] = True
    return node

def to_string(node):
    return f"simplified({node['raw']})"

# Build the full pipeline with pipe()
process = pipe(parse, simplify, to_string)

result = process("x^2 + 2*x + 1")
print("Result:", result)

# compose goes right-to-left
process2 = compose(to_string, simplify, parse)
result2 = process2("x^2 + 2*x + 1")
print("Result2:", result2)
\end{lstlisting}

\pythontutorbox

\textbf{Code Explanation:}

\texttt{pipe} and \texttt{compose} are \textbf{function composition} utilities that chain multiple functions into a single transformation pipeline.

\textbf{What happens here?}
\begin{itemize}
    \item \texttt{pipe(parse, simplify, to\_string)} creates one function that runs all three in order.
    \item \texttt{compose(f, g, h)(x)} is equivalent to \texttt{f(g(h(x)))} --- right to left.
    \item Both use \texttt{functools.reduce} internally to fold the function list.
    \item In the real analyzer: \texttt{pipe(parse, simplify, to\_string)("2*x+1")} gives a clean one-line transformation.
\end{itemize}

\newpage
% ════════════════════════════════════════════════════════════════════════════
%  9. FIXED-POINT COMBINATOR
% ════════════════════════════════════════════════════════════════════════════
\section{9.\quad Fixed-Point Combinator}

\textbf{Code:}

\begin{lstlisting}
# fixed_point: apply a function until output == input
def fixed_point(fn, x, max_iter=100):
    """Apply fn repeatedly until the output equals the input."""
    for step in range(max_iter):
        result = fn(x)
        print(f"  step {step}: {x!r} -> {result!r}")
        if result == x:
            print(f"  Converged after {step} steps.")
            return result
        x = result
    raise RuntimeError("Did not converge")

# Demo 1: integer halving converges to 0
print("=== x // 2 ===")
fixed_point(lambda x: x // 2, 100)

# Demo 2: algebraic simplification rule
# simplify_once: applies a single pass of rules
def simplify_once(expr):
    """Remove 0+ and *1 (toy version)."""
    if expr.startswith("0+"):   return expr[2:]
    if expr.endswith("*1"):     return expr[:-2]
    return expr

print("\n=== simplify_once on '0+x*1' ===")
result = fixed_point(simplify_once, "0+x*1")
print("Final:", result)
\end{lstlisting}

\pythontutorbox

\textbf{Code Explanation:}

The \textbf{fixed-point combinator} applies a function repeatedly until its output is identical to its input. This is how the simplifier reaches a stable, fully simplified form.

\textbf{What happens here?}
\begin{itemize}
    \item \texttt{fixed\_point(fn, x)} calls \texttt{fn(x)}, then \texttt{fn(fn(x))}, and so on.
    \item It stops as soon as one pass produces no change (\texttt{result == x}).
    \item Applied to \texttt{simplify\_once}, the AST is simplified until no more rules fire.
    \item This is a clean, declarative alternative to writing a loop that tracks ``did anything change?''
\end{itemize}

\newpage
% ════════════════════════════════════════════════════════════════════════════
%  10. RECURSION
% ════════════════════════════════════════════════════════════════════════════
\section{10.\quad Recursion}

\textbf{Code:}

\begin{lstlisting}
from collections import namedtuple

Number   = namedtuple("Number",   ["value"])
Variable = namedtuple("Variable", ["name"])
BinOp    = namedtuple("BinOp",    ["op", "left", "right"])
UnaryOp  = namedtuple("UnaryOp",  ["op", "operand"])

def differentiate(node, var):
    """Recursively differentiate an AST node w.r.t. var."""
    t = type(node)

    if t is Number:                      # base case
        return Number(0)

    if t is Variable:                    # base case
        return Number(1) if node.name == var else Number(0)

    if t is UnaryOp:                     # recursive case
        return UnaryOp("-", differentiate(node.operand, var))

    if t is BinOp:                       # recursive case
        dl = differentiate(node.left,  var)
        dr = differentiate(node.right, var)
        if node.op == "+":
            return BinOp("+", dl, dr)   # sum rule
        if node.op == "*":              # product rule
            return BinOp("+",
                BinOp("*", dl, node.right),
                BinOp("*", node.left,  dr))

# Differentiate x*x w.r.t. x
expr = BinOp("*", Variable("x"), Variable("x"))
deriv = differentiate(expr, "x")
print("d/dx(x*x) =", deriv)
\end{lstlisting}

\pythontutorbox

\textbf{Code Explanation:}

\textbf{Recursion} is when a function calls itself. The differentiator recurses into sub-expressions (left and right children) until it reaches a base case (\texttt{Number} or \texttt{Variable}).

\textbf{What happens here?}
\begin{itemize}
    \item \texttt{differentiate(BinOp("*", x, x), "x")} calls itself twice to get \texttt{dl} and \texttt{dr}.
    \item Each recursive call processes a simpler sub-tree (\texttt{Variable("x")}).
    \item The base cases (\texttt{Number} $\to$ 0, same \texttt{Variable} $\to$ 1) stop the recursion.
    \item The product rule \(d(uv)/dx = u'v + uv'\) is assembled from the returned values.
\end{itemize}

\newpage
% ════════════════════════════════════════════════════════════════════════════
%  11. FILTER
% ════════════════════════════════════════════════════════════════════════════
\section{11.\quad \texttt{ffilter()} --- Curried Functional Filter}

\textbf{Code:}

\begin{lstlisting}
# ffilter: a @curry-decorated functional wrapper around filter()
# Source: equation_parser/utils.py
from functools import reduce

def curry(fn):
    def _outer(a):
        def _inner(b):
            return fn(a, b)
        return _inner
    return _outer

@curry
def ffilter(predicate, iterable):
    """Curried functional filter over an iterable."""
    return filter(predicate, iterable)

# @curry lets us call it in two ways:

# Way 1: partial application -- supply the predicate first
is_even = lambda x: x % 2 == 0
filter_evens = ffilter(is_even)      # returns a function
result = list(filter_evens([1, 2, 3, 4, 5, 6]))
print("Evens:", result)              # [2, 4, 6]

# Way 2: use directly in pipe() without a lambda wrapper
from functools import reduce as fred
def pipe(*fns):
    return lambda x: fred(lambda a, f: f(a), fns, x)

from collections import namedtuple
Variable = namedtuple("Variable", ["name"])
Number   = namedtuple("Number",   ["value"])
BinOp    = namedtuple("BinOp",    ["op", "left", "right"])

# Filter: keep only nodes with value > 10
no_small = ffilter(lambda n: type(n) is Number and n.value > 10)
nodes = [Number(5), Number(15), Number(3), Number(20)]
print("Numbers > 10:", list(no_small(nodes)))
\end{lstlisting}

\pythontutorbox

\textbf{Code Explanation:}

\texttt{ffilter} is a \textbf{curried} wrapper around Python's built-in \texttt{filter()}. The \texttt{@curry} decorator lets it be applied one argument at a time.

\textbf{What happens here?}
\begin{itemize}
    \item \texttt{ffilter(predicate)} returns a \textbf{new function} that waits for the iterable.
    \item This means it can be placed inside a \texttt{pipe()} without wrapping in a \texttt{lambda}.
    \item \texttt{filter\_evens = ffilter(is\_even)} creates a reusable even-number filter.
    \item Curried functions compose cleanly: \texttt{pipe(ffilter(is\_even), fmap(double), freduce(add))}.
\end{itemize}


\textbf{Code:}

\begin{lstlisting}
# filter(): select AST nodes that satisfy a condition
from collections import namedtuple

Number   = namedtuple("Number",   ["value"])
Variable = namedtuple("Variable", ["name"])
BinOp    = namedtuple("BinOp",    ["op", "left", "right"])

# Flatten an AST into a list of all nodes (generator)
def all_nodes(node):
    yield node
    if type(node) is BinOp:
        yield from all_nodes(node.left)
        yield from all_nodes(node.right)

# Build expression: x + 2 + y + 5
expr = BinOp("+", BinOp("+", Variable("x"), Number(2)),
                  BinOp("+", Variable("y"), Number(5)))

nodes = list(all_nodes(expr))
print("All nodes:", nodes)

# filter: keep only Variable nodes
variables = list(filter(lambda n: type(n) is Variable, nodes))
print("Variables:", variables)

# filter: keep only Number nodes
numbers = list(filter(lambda n: type(n) is Number, nodes))
print("Numbers:", numbers)

# filter: keep only numbers greater than 3
big = list(filter(lambda n: type(n) is Number and n.value > 3, nodes))
print("Numbers > 3:", big)
\end{lstlisting}

\pythontutorbox

\textbf{Code Explanation:}

\texttt{filter()} selects elements from an iterable that satisfy a predicate function.

\textbf{What happens here?}
\begin{itemize}
    \item \texttt{all\_nodes()} is a generator that yields every node in the AST.
    \item \texttt{filter(lambda n: type(n) is Variable, nodes)} keeps only \texttt{Variable} nodes.
    \item The predicate \texttt{lambda n: type(n) is Number and n.value > 3} filters on multiple conditions.
    \item In the real analyzer, \texttt{filter} is used to select available operations or valid sub-expressions.
\end{itemize}

\newpage
% ════════════════════════════════════════════════════════════════════════════
%  12. MAP
% ════════════════════════════════════════════════════════════════════════════
\section{12.\quad \texttt{fmap()} --- Curried Functional Map}

\textbf{Code:}

\begin{lstlisting}
# fmap: a @curry-decorated functional wrapper around map()
# Source: equation_parser/utils.py
from functools import reduce

def curry(fn):
    def _outer(a):
        def _inner(b):
            return fn(a, b)
        return _inner
    return _outer

@curry
def fmap(fn, iterable):
    """Curried functional map over an iterable."""
    return map(fn, iterable)

# Partial application: fix the transform function first
double = lambda x: x * 2
doubler = fmap(double)          # returns a function waiting for an iterable
result = list(doubler([1, 2, 3]))
print("Doubled:", result)       # [2, 4, 6]

# Use with AST nodes from the equation parser
from collections import namedtuple
Number = namedtuple("Number", ["value"])

nodes = [Number(1), Number(2), Number(3)]

# Map: extract all values
values = list(fmap(lambda n: n.value)(nodes))
print("Values:", values)   # [1, 2, 3]

# Map: negate all Number values (returns new immutable nodes)
neg = list(fmap(lambda n: Number(-n.value))(nodes))
print("Negated:", neg)     # [Number(-1), Number(-2), Number(-3)]
\end{lstlisting}

\pythontutorbox

\textbf{Code Explanation:}

\texttt{fmap} is a \textbf{curried} wrapper around Python's built-in \texttt{map()}.

\textbf{What happens here?}
\begin{itemize}
    \item \texttt{fmap(double)} returns a function that maps \texttt{double} over any iterable.
    \item Because \texttt{fmap} is curried, it plugs seamlessly into \texttt{pipe()} pipelines.
    \item \texttt{fmap(lambda n: n.value)(nodes)} applies a transform to every AST node.
    \item All transformations produce \textbf{new} objects --- immutability is preserved.
\end{itemize}


\textbf{Code:}

\begin{lstlisting}
from collections import namedtuple

Number   = namedtuple("Number",   ["value"])
Variable = namedtuple("Variable", ["name"])
BinOp    = namedtuple("BinOp",    ["op", "left", "right"])

# to_string: convert any node to a readable string
def to_string(node):
    t = type(node)
    if t is Number:   return str(node.value)
    if t is Variable: return node.name
    if t is BinOp:
        return f"({to_string(node.left)} {node.op} {to_string(node.right)})"

# A list of AST nodes
nodes = [
    Number(42),
    Variable("x"),
    BinOp("+", Number(1), Variable("y")),
    BinOp("*", Variable("a"), Number(3)),
]

# map(): apply to_string to every node
strings = list(map(to_string, nodes))
print("Strings:", strings)

# map(): extract the type name of each node
type_names = list(map(lambda n: type(n).__name__, nodes))
print("Types:", type_names)

# map(): negate all Number values
numbers = [Number(1), Number(2), Number(3)]
negated = list(map(lambda n: Number(-n.value), numbers))
print("Negated:", negated)
\end{lstlisting}

\pythontutorbox

\textbf{Code Explanation:}

\texttt{map()} applies a function to \textbf{every element} of an iterable and returns the transformed results.

\textbf{What happens here?}
\begin{itemize}
    \item \texttt{map(to\_string, nodes)} converts each AST node into its string representation.
    \item \texttt{map(lambda n: type(n).\_\_name\_\_, nodes)} extracts the class name of each node.
    \item \texttt{map(lambda n: Number(-n.value), numbers)} creates new negated \texttt{Number} nodes.
    \item In the real analyzer, \texttt{map} is used to transform collections of nodes uniformly.
\end{itemize}

\newpage
% ════════════════════════════════════════════════════════════════════════════
%  13. REDUCE
% ════════════════════════════════════════════════════════════════════════════
\section{13.\quad \texttt{freduce()} --- Curried Functional Reduce}

\textbf{Code:}

\begin{lstlisting}
# freduce: a @curry-decorated wrapper around functools.reduce()
# Source: equation_parser/utils.py
from functools import reduce

def curry(fn):
    def _outer(a):
        def _inner(b):
            return fn(a, b)
        return _inner
    return _outer

@curry
def freduce(fn, iterable):
    """Curried functional reduce over an iterable."""
    return reduce(fn, iterable)

# Partial application: fix the combining function first
add = lambda a, b: a + b
summer = freduce(add)        # returns a function waiting for an iterable
result = summer([1, 2, 3, 4])
print("Sum:", result)        # 10  (1+2+3+4)

# pipe: filter evens -> double -> sum (all curried, no lambdas needed)
@curry
def ffilter(predicate, iterable): return filter(predicate, iterable)
@curry
def fmap(fn, iterable): return map(fn, iterable)

def pipe(*fns):
    return lambda x: reduce(lambda a, f: f(a), fns, x)

sum_doubled_evens = pipe(
    ffilter(lambda x: x % 2 == 0),  # keep evens
    fmap(lambda x: x * 2),           # double each
    freduce(lambda a, b: a + b),     # sum all
)
print("Result:", sum_doubled_evens([1, 2, 3, 4, 5]))
# [2,4] -> [4,8] -> 12
\end{lstlisting}

\pythontutorbox

\textbf{Code Explanation:}

\texttt{freduce} is a \textbf{curried} wrapper around \texttt{functools.reduce()}. It reduces an iterable to a single value by repeatedly applying a combining function.

\textbf{What happens here?}
\begin{itemize}
    \item \texttt{freduce(add)} returns a summer function waiting for a list.
    \item The full \texttt{pipe(ffilter(...), fmap(...), freduce(...))} pipeline needs \textbf{no lambdas} because all three are curried.
    \item \texttt{pipe} uses \texttt{reduce} internally to pass the value through the chain.
    \item This is the cleanest form of the map-filter-reduce pattern in functional Python.
\end{itemize}


\textbf{Code:}

\begin{lstlisting}
from functools import reduce
from collections import namedtuple

Number = namedtuple("Number", ["value"])

# reduce(): fold a list of functions left-to-right into one call
# This is exactly how pipe() works internally

def pipe(*fns):
    """Left-to-right function composition using reduce."""
    def _piped(x):
        return reduce(lambda acc, f: f(acc), fns, x)
    return _piped

# Individual transformation steps
double     = lambda x: x * 2
add_ten    = lambda x: x + 10
to_string  = lambda x: f"result={x}"

# Manually trace what reduce does:
fns = [double, add_ten, to_string]
x = 5
print("Start:", x)
step1 = double(x);         print("After double:", step1)
step2 = add_ten(step1);    print("After add_ten:", step2)
step3 = to_string(step2);  print("After to_string:", step3)

# pipe does the same automatically via reduce
pipeline = pipe(double, add_ten, to_string)
print("\npipe result:", pipeline(5))

# reduce to add a list of Numbers
nums = [Number(10), Number(20), Number(30)]
total = reduce(lambda a, b: Number(a.value + b.value), nums)
print("Sum of Number nodes:", total)
\end{lstlisting}

\pythontutorbox

\textbf{Code Explanation:}

\texttt{reduce()} applies a function cumulatively to sequence elements, reducing them to a single value.

\textbf{What happens here?}
\begin{itemize}
    \item \texttt{pipe()} internally uses \texttt{reduce(lambda acc, f: f(acc), fns, x)} to thread a value through a list of functions.
    \item \texttt{reduce(lambda acc, f: f(acc), [double, add\_ten, to\_string], 5)} first calls \texttt{double(5)}, then \texttt{add\_ten(10)}, then \texttt{to\_string(20)}.
    \item \texttt{reduce(lambda a, b: Number(a.value + b.value), nums)} sums a list of \texttt{Number} nodes.
\end{itemize}

\newpage
% ════════════════════════════════════════════════════════════════════════════
%  14. LIST COMPREHENSION
% ════════════════════════════════════════════════════════════════════════════
\section{14.\quad List Comprehension}

\textbf{Code:}

\begin{lstlisting}
from collections import namedtuple

Number   = namedtuple("Number",   ["value"])
Variable = namedtuple("Variable", ["name"])
BinOp    = namedtuple("BinOp",    ["op", "left", "right"])

# Flatten AST into list of all sub-nodes
def all_nodes(node):
    yield node
    if type(node) is BinOp:
        yield from all_nodes(node.left)
        yield from all_nodes(node.right)

# AST for: x + 2 * y + 5
expr = BinOp("+",
    BinOp("+", Variable("x"), BinOp("*", Number(2), Variable("y"))),
    Number(5))

nodes = list(all_nodes(expr))

# 1. All variable names using list comprehension
var_names = [n.name for n in nodes if type(n) is Variable]
print("Variable names:", var_names)

# 2. All number values > 1
big_numbers = [n.value for n in nodes if type(n) is Number and n.value > 1]
print("Numbers > 1:", big_numbers)

# 3. All operators from BinOp nodes
operators = [n.op for n in nodes if type(n) is BinOp]
print("Operators:", operators)

# 4. Convert all Numbers to their negatives
negated = [Number(-n.value) for n in nodes if type(n) is Number]
print("Negated numbers:", negated)
\end{lstlisting}

\pythontutorbox

\textbf{Code Explanation:}

\textbf{List comprehension} creates a new list in a single concise line using a condition.

\textbf{What happens here?}
\begin{itemize}
    \item \texttt{[n.name for n in nodes if type(n) is Variable]} collects all variable names in the AST.
    \item \texttt{[n.op for n in nodes if type(n) is BinOp]} extracts operators.
    \item Each comprehension applies both a \textit{transformation} (the expression before \texttt{for}) and a \textit{filter} (the \texttt{if} clause).
\end{itemize}

\newpage
% ════════════════════════════════════════════════════════════════════════════
%  15. DICTIONARY COMPREHENSION
% ════════════════════════════════════════════════════════════════════════════
\section{15.\quad Dictionary Comprehension}

\textbf{Code:}

\begin{lstlisting}
from collections import namedtuple

# Symbolic differentiation rules for built-in math functions
# Each rule is: function_name -> lambda(inner) -> derivative_AST
FuncCall = namedtuple("FuncCall", ["name", "arg"])
BinOp    = namedtuple("BinOp",    ["op", "left", "right"])
Number   = namedtuple("Number",   ["value"])
UnaryOp  = namedtuple("UnaryOp",  ["op", "operand"])

# Base rules list: (name, derivative_lambda)
BASE_RULES = [
    ("sin",  lambda u: FuncCall("cos", u)),
    ("cos",  lambda u: UnaryOp("-", FuncCall("sin", u))),
    ("exp",  lambda u: FuncCall("exp", u)),
    ("ln",   lambda u: BinOp("/", Number(1), u)),
    ("sqrt", lambda u: BinOp("/", Number(1), BinOp("*", Number(2), FuncCall("sqrt", u)))),
]

# Dictionary comprehension: build dispatch table from list
DIFF_RULES = {name: rule for name, rule in BASE_RULES}

print("Registered rules:", list(DIFF_RULES.keys()))

# Use a rule
u = FuncCall("x", None)
sin_deriv = DIFF_RULES["sin"](FuncCall("x", None))
print("d/du sin(u) =", sin_deriv)

ln_deriv = DIFF_RULES["ln"](Number(5))
print("d/du ln(5) =", ln_deriv)
\end{lstlisting}

\pythontutorbox

\textbf{Code Explanation:}

\textbf{Dictionary comprehension} builds a dictionary in a single concise line from an iterable.

\textbf{What happens here?}
\begin{itemize}
    \item \texttt{BASE\_RULES} is a list of \texttt{(name, lambda)} pairs.
    \item \texttt{\{name: rule for name, rule in BASE\_RULES\}} converts it into a dispatch dictionary.
    \item The resulting \texttt{DIFF\_RULES} maps each function name to its derivative rule.
    \item This is more compact than writing out each key assignment manually.
\end{itemize}

\newpage
% ════════════════════════════════════════════════════════════════════════════
%  16. ENUMERATE
% ════════════════════════════════════════════════════════════════════════════
\section{16.\quad \texttt{enumerate()}}

\textbf{Code:}

\begin{lstlisting}
from collections import namedtuple

Number   = namedtuple("Number",   ["value"])
Variable = namedtuple("Variable", ["name"])
BinOp    = namedtuple("BinOp",    ["op", "left", "right"])

def all_nodes(node):
    yield node
    if type(node) is BinOp:
        yield from all_nodes(node.left)
        yield from all_nodes(node.right)

# AST for: x + 2 + y
expr = BinOp("+", BinOp("+", Variable("x"), Number(2)), Variable("y"))
nodes = list(all_nodes(expr))

# enumerate() gives index + value while looping
print("All nodes in the AST:")
for i, node in enumerate(nodes, start=1):
    print(f"  [{i}] {type(node).__name__:10s} -> {node}")

# Use enumerate to number the tokens in an expression
def tokenize_simple(text):
    for ch in text.split():
        yield ch

print("\nTokens with indices:")
for idx, tok in enumerate(tokenize_simple("x + 2 * y")):
    print(f"  Token {idx}: {tok!r}")
\end{lstlisting}

\pythontutorbox

\textbf{Code Explanation:}

\texttt{enumerate()} gives both the \textbf{index} and the \textbf{value} while iterating over a sequence.

\textbf{What happens here?}
\begin{itemize}
    \item \texttt{enumerate(nodes, start=1)} pairs each node with a number starting from 1.
    \item This is used to number AST nodes for display or debugging.
    \item \texttt{enumerate(tokenize\_simple(...))} also works on generators --- it is lazy.
    \item Avoids creating a separate counter variable.
\end{itemize}

\newpage
% ════════════════════════════════════════════════════════════════════════════
%  17. ZIP
% ════════════════════════════════════════════════════════════════════════════
\section{17.\quad \texttt{zip()} Function}

\textbf{Code:}

\begin{lstlisting}
# zip(): combine variable names with their values for evaluation
variables = ["x", "y", "z"]
values    = [5.0, 3.0, 7.0]

# zip pairs them element-by-element
pairs = list(zip(variables, values))
print("Paired:", pairs)

# Build the environment dict from two parallel lists
env = dict(zip(variables, values))
print("Environment:", env)

# This is exactly how an evaluator receives variable bindings
from collections import namedtuple
Variable = namedtuple("Variable", ["name"])
Number   = namedtuple("Number",   ["value"])
BinOp    = namedtuple("BinOp",    ["op", "left", "right"])

def evaluate(node, env):
    t = type(node)
    if t is Number:   return float(node.value)
    if t is Variable: return env[node.name]
    if t is BinOp:
        ops = {"+": lambda a,b: a+b, "*": lambda a,b: a*b}
        return ops[node.op](evaluate(node.left, env), evaluate(node.right, env))

# Evaluate x + y given env
expr = BinOp("+", Variable("x"), Variable("y"))
result = evaluate(expr, env)
print(f"x + y = {result}")

# zip to display variable=value pairs
for var, val in zip(variables, values):
    print(f"  {var} = {val}")
\end{lstlisting}

\pythontutorbox

\textbf{Code Explanation:}

\texttt{zip()} combines two or more iterables into pairs (tuples), element by element.

\textbf{What happens here?}
\begin{itemize}
    \item \texttt{zip(variables, values)} pairs each variable name with its numeric value.
    \item \texttt{dict(zip(variables, values))} builds the environment dictionary directly.
    \item The environment \texttt{env} is then passed to \texttt{evaluate()} to resolve variable nodes.
\end{itemize}

\newpage
% ════════════════════════════════════════════════════════════════════════════
%  18. PARTIAL APPLICATION
% ════════════════════════════════════════════════════════════════════════════
\section{18.\quad Partial Application}

\textbf{Code:}

\begin{lstlisting}
from functools import partial
from collections import namedtuple

Number   = namedtuple("Number",   ["value"])
Variable = namedtuple("Variable", ["name"])
BinOp    = namedtuple("BinOp",    ["op", "left", "right"])
UnaryOp  = namedtuple("UnaryOp",  ["op", "operand"])

def differentiate(node, var):
    t = type(node)
    if t is Number:   return Number(0)
    if t is Variable: return Number(1) if node.name == var else Number(0)
    if t is BinOp:
        dl = differentiate(node.left, var)
        dr = differentiate(node.right, var)
        if node.op == "+": return BinOp("+", dl, dr)
        if node.op == "*":
            return BinOp("+", BinOp("*", dl, node.right),
                              BinOp("*", node.left, dr))
    return node

# Partial application: fix the variable argument
diff_x = partial(differentiate, var="x")   # always differentiates w.r.t. x
diff_y = partial(differentiate, var="y")   # always differentiates w.r.t. y

expr = BinOp("+", Variable("x"), Variable("y"))

print("d/dx (x + y) =", diff_x(expr))
print("d/dy (x + y) =", diff_y(expr))

# Use as a pipeline step
from functools import reduce
def pipe(*fns):
    return lambda x: reduce(lambda a, f: f(a), fns, x)

def to_str(n):
    t = type(n)
    if t is Number: return str(n.value)
    if t is Variable: return n.name
    if t is BinOp: return f"({to_str(n.left)} {n.op} {to_str(n.right)})"

# pipe: parse -> diff_x -> to_str
pipeline = pipe(diff_x, to_str)
print("Pipeline result:", pipeline(expr))
\end{lstlisting}

\pythontutorbox

\textbf{Code Explanation:}

\textbf{Partial application} fixes one or more arguments of a function, producing a new function with fewer arguments.

\textbf{What happens here?}
\begin{itemize}
    \item \texttt{partial(differentiate, var="x")} creates \texttt{diff\_x} --- a version of \texttt{differentiate} with \texttt{var} pre-filled to \texttt{"x"}.
    \item \texttt{diff\_x(expr)} only needs the node; the variable is already fixed.
    \item This makes \texttt{diff\_x} directly usable inside a \texttt{pipe()} without extra \texttt{lambda} wrappers.
\end{itemize}

\newpage
% ════════════════════════════════════════════════════════════════════════════
%  19. LAMBDA
% ════════════════════════════════════════════════════════════════════════════
\section{19.\quad Lambda (Anonymous Functions)}

\textbf{Code:}

\begin{lstlisting}
from collections import namedtuple

Number = namedtuple("Number", ["value"])
Variable = namedtuple("Variable", ["name"])
BinOp = namedtuple("BinOp", ["op", "left", "right"])

# Lambda functions as VALUES in the operator dispatch table
BIN_OPS = {
    "+": lambda a, b: a + b,      # anonymous add
    "-": lambda a, b: a - b,      # anonymous subtract
    "*": lambda a, b: a * b,      # anonymous multiply
    "/": lambda a, b: a / b,      # anonymous divide
    "^": lambda a, b: a ** b,     # anonymous power
}

# Evaluate using the lambda dispatch table
def evaluate(node, env):
    t = type(node)
    if t is Number:   return float(node.value)
    if t is Variable: return float(env[node.name])
    if t is BinOp:
        l = evaluate(node.left,  env)
        r = evaluate(node.right, env)
        return BIN_OPS[node.op](l, r)

# Test each operator
for op in ["+", "-", "*", "/", "^"]:
    node = BinOp(op, Number(10), Number(2))
    result = evaluate(node, {})
    print(f"10 {op} 2 = {result}")

# Lambda in sorted()
items = [("x", 5), ("a", 2), ("m", 9)]
by_value = sorted(items, key=lambda pair: pair[1])
print("Sorted by value:", by_value)
\end{lstlisting}

\pythontutorbox

\textbf{Code Explanation:}

A \textbf{lambda} is an anonymous (unnamed) function written in a single expression. It is used wherever a short function is needed as a value.

\textbf{What happens here?}
\begin{itemize}
    \item \texttt{lambda a, b: a + b} creates an anonymous addition function stored in \texttt{BIN\_OPS["+"]}.
    \item All five operators are expressed as lambdas, eliminating repetitive \texttt{if/else} branches.
    \item \texttt{sorted(..., key=lambda pair: pair[1])} uses a lambda as a sort key.
\end{itemize}

\newpage
% ════════════════════════════════════════════════════════════════════════════
%  20. FULL PIPELINE
% ════════════════════════════════════════════════════════════════════════════
\section{20.\quad Full Functional Pipeline}

\textbf{Code:}

\begin{lstlisting}
from collections import namedtuple
from functools import reduce

# ── AST Nodes ──
Number   = namedtuple("Number",   ["value"])
Variable = namedtuple("Variable", ["name"])
BinOp    = namedtuple("BinOp",    ["op", "left", "right"])
Token    = namedtuple("Token",    ["kind", "value"])
EOF = "EOF"

def pipe(*fns):
    return lambda x: reduce(lambda a, f: f(a), fns, x)

# ── Generator tokenizer ──
def tokenize(text):
    pos = 0
    while pos < len(text):
        ch = text[pos]
        if ch.isspace():  pos += 1; continue
        if ch.isdigit():
            s = pos
            while pos < len(text) and text[pos].isdigit(): pos += 1
            yield Token("NUM", int(text[s:pos])); continue
        if ch.isalpha(): yield Token("VAR", ch); pos += 1; continue
        if ch in "+-*/": yield Token("OP", ch); pos += 1; continue
        pos += 1
    yield Token(EOF, None)

# ── Parser ──
def parse(text):
    toks = list(tokenize(text))
    pos = [0]
    def cur():  return toks[pos[0]]
    def eat():
        t = toks[pos[0]]; pos[0] += 1; return t
    def atom():
        t = eat()
        if t.kind == "NUM": return Number(t.value)
        if t.kind == "VAR": return Variable(t.value)
    def expr():
        left = atom()
        while cur().kind == "OP":
            op = eat().value; left = BinOp(op, left, atom())
        return left
    return expr()

# ── Simplifier with fixed_point ──
def simp(n):
    if type(n) is BinOp:
        l, r = simp(n.left), simp(n.right)
        if type(l) is Number and type(r) is Number:
            ops = {"+": l.value+r.value, "*": l.value*r.value,
                   "-": l.value-r.value}
            return Number(ops[n.op])
        return BinOp(n.op, l, r)
    return n

def simplify(n):
    for _ in range(50):
        r = simp(n)
        if r == n: return r
        n = r
    return n

# ── Pretty-printer ──
def to_str(n):
    if type(n) is Number:   return str(n.value)
    if type(n) is Variable: return n.name
    if type(n) is BinOp:    return f"({to_str(n.left)} {n.op} {to_str(n.right)})"

# ── FULL PIPELINE ──
process = pipe(parse, simplify, to_str)

for expr in ["x + 2 + 3", "2 * 3 + y", "1 + 2 + 3"]:
    print(f"  {expr:20s}  ->  {process(expr)}")
\end{lstlisting}

\pythontutorbox

\textbf{Code Explanation:}

The \textbf{full pipeline} chains every functional concept together using \texttt{pipe()}: generator tokenizer $\to$ parser $\to$ fixed-point simplifier $\to$ pretty-printer.

\textbf{What happens here?}
\begin{itemize}
    \item \texttt{pipe(parse, simplify, to\_str)} creates a single transformation function.
    \item Each step receives the \textbf{output} of the previous step as its input.
    \item The string \texttt{"2 * 3 + y"} flows through: tokenize $\to$ AST $\to$ simplified AST $\to$ \texttt{"(6 + y)"}.
    \item This demonstrates how all concepts --- generators, iterators, higher-order functions, recursion, fixed-point, lambdas, and dispatch tables --- compose into one clean system.
\end{itemize}



\newpage
% ════════════════════════════════════════════════════════════════════════════
%  21. PLOTTER MODULE
% ════════════════════════════════════════════════════════════════════════════
\section{21.\quad Plotter --- Visualizing Equations with Matplotlib}

\textbf{Code:}

\begin{lstlisting}
# Source: equation_parser/plotter.py
import numpy as np
import matplotlib.pyplot as plt

# Assume parse, evaluate, differentiate, simplify, to_string, variable_set
# are already available from the equation_parser package.

def _safe_eval(tree, var, x_val):
    """Evaluate tree at a single point; return nan on math errors."""
    try:
        return evaluate(tree, {var: float(x_val)})
    except (ValueError, ZeroDivisionError, OverflowError, NameError):
        return float('nan')

def _vectorized_eval(tree, var, xs):
    """Evaluate tree over a NumPy array, handling domain errors."""
    return [_safe_eval(tree, var, x) for x in xs]

def plot_expression(expr_str, x_range=(-10, 10), num_points=500):
    """Plot a single expression string."""
    tree = parse(expr_str)
    vs = variable_set(tree)
    var = next(iter(vs)) if vs else "x"
    xs = list(range(int(x_range[0]), int(x_range[1]) + 1))
    ys = _vectorized_eval(tree, var, xs)
    print(f"Plotting f({var}) = {expr_str}")
    for x, y in zip(xs, ys):
        print(f"  {var}={x:4}  f={y}")

def plot_with_derivative(expr_str, x_range=(-5, 5)):
    """Plot expression and its symbolic derivative."""
    tree = parse(expr_str)
    vs = variable_set(tree)
    var = next(iter(vs)) if vs else "x"
    d_tree = simplify(differentiate(tree, var))
    d_str  = to_string(d_tree)
    xs = list(range(int(x_range[0]), int(x_range[1]) + 1))
    ys  = _vectorized_eval(tree, var, xs)
    dys = _vectorized_eval(d_tree, var, xs)
    print(f"f ({var}) = {expr_str}")
    print(f"f'({var}) = {d_str}")
    for x, y, dy in zip(xs, ys, dys):
        print(f"  {var}={x:3}  f={y:8.3f}  f'={dy:8.3f}")
\end{lstlisting}

\pythontutorbox

\textbf{Code Explanation:}

The \textbf{plotter} module integrates the parser, evaluator, and differentiator with Matplotlib. All functions are \textbf{pure} --- they take expressions as strings and return figures.

\textbf{What happens here?}
\begin{itemize}
    \item \texttt{\_safe\_eval} evaluates one point, returning \texttt{nan} on errors (handles \texttt{log(0)}, \texttt{1/0}, etc.).
    \item \texttt{\_vectorized\_eval} maps \texttt{\_safe\_eval} over a NumPy array --- a \textbf{higher-order function} applied across all x-values.
    \item \texttt{plot\_expression} parses the expression, auto-detects the variable, samples it, and draws the plot.
    \item \texttt{plot\_with\_derivative} differentiates symbolically first (using \texttt{differentiate + simplify}), then plots both curves.
    \item REPL commands: \texttt{:plot x\char`\^2}, \texttt{:plotm x\char`\^2, sin(x)}, \texttt{:plotd x\char`\^3 [-5:5]}.
\end{itemize}

\newpage
% ════════════════════════════════════════════════════════════════════════════
%  22. UPDATED FULL PIPELINE WITH fmap / ffilter / freduce
% ════════════════════════════════════════════════════════════════════════════
\section{22.\quad Updated Full Pipeline --- fmap, ffilter, freduce}

\textbf{Code:}

\begin{lstlisting}
# The updated pipeline now uses the REAL curried functions from utils.py
# Source: equation_parser/utils.py  +  tests/test_equation.py
from functools import reduce

# --- curry ---
def curry(fn):
    return lambda a: lambda b: fn(a, b)

# --- pipe ---
def pipe(*fns):
    return lambda x: reduce(lambda a, f: f(a), fns, x)

# --- curried functional trio ---
@curry
def fmap(fn, iterable):       return map(fn, iterable)

@curry
def ffilter(pred, iterable):  return filter(pred, iterable)

@curry
def freduce(fn, iterable):    return reduce(fn, iterable)

# ═══ DEMO 1: sum of doubled even numbers ════════════════
add    = lambda a, b: a + b
is_even = lambda x: x % 2 == 0
double  = lambda x: x * 2

sum_of_doubled_evens = pipe(
    ffilter(is_even),        # [1,2,3,4,5] -> filter -> [2,4]
    fmap(double),            # [2,4]        -> map    -> [4,8]
    freduce(add),            # [4,8]        -> reduce -> 12
)
result = sum_of_doubled_evens([1, 2, 3, 4, 5])
print("Sum of doubled evens:", result)  # 12

# ═══ DEMO 2: full equation pipeline ════════════════════
from collections import namedtuple
Number   = namedtuple("Number",   ["value"])
Variable = namedtuple("Variable", ["name"])
BinOp    = namedtuple("BinOp",    ["op", "left", "right"])

def to_str(n):
    t = type(n)
    if t is Number:   return str(n.value)
    if t is Variable: return n.name
    if t is BinOp:    return f"({to_str(n.left)} {n.op} {to_str(n.right)})"

# collect all values from a list of Number nodes
extract_values = pipe(
    ffilter(lambda n: type(n) is Number),
    fmap(lambda n: n.value),
    freduce(lambda a, b: a + b),
)
nodes = [BinOp("+", Number(10), Number(5)), Number(3), Number(7)]
print("Sum of Number values in list:", extract_values(nodes))  # 10
\end{lstlisting}

\pythontutorbox

\textbf{Code Explanation:}

This section shows the \textbf{real, updated} pipeline from the codebase, combining all functional tools: \texttt{curry}, \texttt{pipe}, \texttt{fmap}, \texttt{ffilter}, and \texttt{freduce}.

\textbf{What happens here?}
\begin{itemize}
    \item \texttt{ffilter(is\_even)} partially applies the predicate --- the result is a function that takes any iterable.
    \item \texttt{fmap(double)} partially applies the transform.
    \item \texttt{freduce(add)} partially applies the combiner.
    \item All three plug into \texttt{pipe()} \textbf{without any extra lambda wrappers} --- this is the power of currying combined with higher-order composition.
    \item \texttt{pipe(ffilter(is\_even), fmap(double), freduce(add))([1,2,3,4,5])} produces \texttt{12}.
\end{itemize}

\end{document}
